%%%%%
%%%%%  Use LUALATEX, not LATEX.
%%%%%
%%%%
\documentclass[]{VUMIFTemplateClass}

\usepackage{indentfirst}
\usepackage{amsmath, amsthm, amssymb, amsfonts}
\usepackage{mathtools}
\usepackage{physics}
\usepackage{graphicx}
\usepackage{verbatim}
\usepackage[hidelinks]{hyperref}
\usepackage{color,algorithm,algorithmic}
\usepackage[nottoc]{tocbibind}
\usepackage{tocloft}

\usepackage{titlesec}
\newcommand{\sectionbreak}{\clearpage}

\makeatletter
\renewcommand{\fnum@algorithm}{\thealgorithm}
\makeatother
\renewcommand\thealgorithm{\arabic{algorithm} algorithm}

\usepackage{biblatex}
\bibliography{bibliografija}
%% to change the numbering (numeric or alphabetic) of bibliographic sources, make the change in VUMIFTemplateClass.cls, line 139

% Author's MACROS
\newcommand{\EE}{\mathbb{E}\,} % Mean
\newcommand{\ee}{{\mathrm e}}  % nice exponent
\newcommand{\RR}{\mathbb{R}}


\studyprogramme{Software engineering} %Write your study programme (example – Software engineering, Financial and Actuarial Mathematics, etc.)
% \worktype{\{Work type\}} % Bachelor's thesis or Master's thesis
\worktitle{Human computer interaction}
\secondworktitle{Žmogaus kompiuterio sąveika}
\workauthor{Tadas Riksas, Darius Spruogis, Gustas Mickus, Kajus Bicka}

%There may be more than one author, in which case each author is written from a new line which is added in Titlepage.tex or LongerTitlePage.tex
%\secondauthor{Name Surname} %If present, otherwise delete

\supervisor{Kristina Lapin}
\reviewer{Kristina Lapin} %If present, otherwise delete
% \scientificadvisor{pedagogical/scientific title Name Surname} %If present, otherwise delete

\begin{document}
\selectlanguage{english}

\onehalfspacing
\begin{titlepage}
\vskip 20pt
\begin{center}
\includegraphics[scale=0.55]{images/MIF.png}
\end{center}

\makeatletter

\vskip 20pt
\centerline{\bf \large \textbf{VILNIUS UNIVERSITY}}
\vskip 10pt
\centerline{\large \textbf{FACULTY OF MATHEMATICS AND INFORMATICS}}
\vskip 10pt
\centerline{\large \textbf{\MakeUppercase{\@studyprogramme \space study programme}}}

\vskip 80pt
\centerline{\Large \@worktype}
\vskip 20pt
\begin{center}
    {\bf \LARGE \@worktitle}
\end{center}
\begin{center}
    {\bf \Large \@secondworktitle}
\end{center}
\vskip 80pt

\centering{\Large \@workauthor}
\@ifundefined{@secondauthor}{}
{
\vskip 10pt
\centering{\Large \@secondauthor}
}
\vskip 20pt

\centering{
    \begin{tabular}{rcp{.7\textwidth}}
        {\Large Supervisor} & {\Large :} & {\Large \@supervisor}\\[10pt]
        \@ifundefined{@scientificadvisor}{}
            {
                {\Large Scientific advisor} & {\Large :} & {\Large \@scientificadvisor}\\[10pt]
            }
        \@ifundefined{@reviewer}{}
            {
                {\Large Reviewer} & {\Large :} & {\Large \@reviewer}\\[10pt]
            }
    \end{tabular}}


\vskip 110pt

\centerline{\large \textbf{Vilnius}}
\centerline{\large \textbf{\the\year{}}}

\makeatother

\newpage
\end{titlepage}
%\newgeometry{top=2cm,bottom=2cm,right=2cm,left=3cm}
\setcounter{page}{2}


\singlespacing
\selectlanguage{english}
% list of figures, delete if not needed
% \listoffigures 

%list of tables, delete if not needed
% \listoftables

%Turinys
\tableofcontents
\onehalfspacing


% %Section for abbreviations
% \sectionnonum{List of abbreviations} %Leave if necessary
% This section is for when abbreviations are used. For example:

% \begin{tabular}{rcp{.7\textwidth}}
%     {u.d.i.r.v.} & {} & {uniformly distributed independent random variables}
% \end{tabular}

\section{Introduction}
\subsection{Project Title}

\textbf{Title:} UniServe

\textbf{Explanation:} "Uni" represents \emph{universal} inidicating the platform's flexibility across a variety of business models 
(reservation based businesses such as massage services or order based businesses such as restaurants). 
"Serve" refers to \emph{service}, emphasizing the platform’s role in providing core functionalities for these types of businesses.

\subsection{Problem statement}

\subsubsection{The background}

% Decades ago, it was still common to go to restaurants and expect menus to be
% handed out on-site, or to call by phone to register an appointment. However, as
% we move into an age of ever-accelerating communication and information exchange,
% these methods have become stubborn and slow. More customers and businesses recognize a need for
% faster communication and more accurate information exchange.

% Reliance on traditional methods slows down information exchange between
% employees, introduces human error, and creates friction in customer
% interactions. This leads to employee dissatisfaction, increased workload, as
% well as a decrease in operational efficiency.

% Slower communications and unstandardized registrations or menus introduce friction
% in customer interactions, makes customers more reluctant to engage with the
% business and leads to overall decrease in customer satisfaction.


% Traditional methods like handing out menus and taking orders on-site or booking appointments by phone have become slow and inefficient in the digital age. Small to medium-sized businesses in restaurant and beauty industries often cannot afford specialized web applications due to limited budgets and technical expertise.
% This creates multiple problems: employees face increased workload and inefficient information exchange, while customers increasingly expect online services like remote ordering and appointment booking. Without accessible digital platforms, businesses risk missing potential customers who rely on the internet, leading to reduced customer reach, lower operational efficiency, and decreased customer satisfaction. The lack of automation results in higher recurring costs and makes customers more reluctant to engage with businesses.


\subsubsection{Impact on the business}

Owners of small to medium-sized businesses in the restaurant and beauty
industries often cannot afford the resources to develop and maintain a
specialized web application for managing their business operations and customer
interactions. This is due to many of these businesses having limited budgets and
technical expertise, which makes custom software solutions expensive.

The people affected by this problem include customers, who increasingly expect
online services such as remote ordering or appointment booking. This challenge
is common in both urban and suburban areas where internet use is high and the
contest for the attention and loyalty of customers has become a defining force.
Without an accessible digital platform, businesses risk missing potential
customers who rely on the internet to discover and engage with services,
leading to reduced customer reach. 

Furthermore, business owners also face operational inefficiencies and lower
profits due to the lack of automation, requiring increased employee workload,
leading to increased recurring costs. The problem occurs continuously throughout
the year, but becomes critical during peak business periods such as holidays,
weekends, or promotional events.

\section{Stakeholder analysis}

\subsection{Identifying stakeholders}


% we have identified our primary stakeholders as employees, secondary as customers, indirect as business owners, competitors as null and development and maintance team as sysadmins and tech support. create a table bellow

\begin{table}[h]
  \centering
  \begin{tabular}{|c|c|}
    \hline
    Stakeholder Type    & Stakeholders \\ \hline
    Primary             & Client Business Owners, Employees \\ \hline
    Secondary           & Customers \\ \hline
    Indirect (Tertiary) & App Business Owners, Managers \\ \hline
    External/Competitors & Square, Toast POS users \\ \hline
    Technical Support   & Sysadmins, Tech Support \\ \hline
  \end{tabular}
  \caption{Stakeholders Identification}
  \label{tab:stakeholders}
\end{table}

\textbf{Primary Stakeholders} directly and frequently use the system. Client
business owners (restaurant/salon owners) operate the system daily. Employees
process orders and manage appointments through the application.

\textbf{Secondary Stakeholders} are occasionally involved or interact via
others. Customers use the system when placing orders or booking appointments.
Managers oversee operations and may need access to system statistics.

\textbf{Indirect (Tertiary) Stakeholders} influence purchase decisions and
benefit from the system. App business owners (our company's owners) profit from
the system's success. Business managers make decisions about adopting the
system.

\textbf{Competitors} represent existing market solutions. Users currently
employing Square, Toast POS, or similar platforms are potential adopters of our
system.

\textbf{Development and Maintenance Team} creates and maintains the system.
Sysadmins handle infrastructure and servers. Tech support staff resolve user
issues and provide assistance.

\subsection{Stakeholder Expectations}

\textbf{Primary Stakeholders:}
\begin{itemize}
    \item \textbf{Client Business Owners} expect comprehensive business management tools including revenue tracking, staff oversight, inventory management, and detailed analytics to improve operational efficiency and reduce costs.
    \item \textbf{Employees} expect an intuitive, fast interface that streamlines order processing and appointment scheduling without complex training requirements.
\end{itemize}

\textbf{Secondary Stakeholders:}
\begin{itemize}
    \item \textbf{Customers} expect a simple, reliable booking and ordering experience with quick response times and accurate information.
    \item \textbf{Managers} expect access to performance statistics, staff productivity metrics, and operational reports for decision-making.
\end{itemize}

\textbf{Indirect (Tertiary) Stakeholders:}
\begin{itemize}
    \item \textbf{App Business Owners} expect sustainable revenue growth, market expansion, and competitive positioning in the business management software sector.
    \item \textbf{Business Managers} expect clear ROI demonstration, cost-effectiveness compared to existing solutions, and smooth implementation processes.
\end{itemize}

\textbf{Competitors:}
\begin{itemize}
    \item \textbf{Square/Toast POS users} expect feature parity or superiority, competitive pricing, and seamless migration from existing systems.
\end{itemize}

\textbf{Development and Maintenance Team:}
\begin{itemize}
    \item \textbf{Sysadmins} expect reliable system architecture, scalable infrastructure, and clear deployment procedures.
    \item \textbf{Tech Support} expect comprehensive documentation, effective troubleshooting tools, and manageable user issue resolution processes.
\end{itemize}






\section{User needs analysis}




% \sectionnonum{Results and conclusions}
% For details of what needs to be written in this section, please refer to the methodology requirements of the respective programme. 

\printbibliography[title = {References and sources}]

\appendix
\renewcommand{\thesection}{Appendix \arabic{section}. }

\section{\phantom{Appendix} Examples of citations}
In the document \textit{bibliography.bib}, you need to add all the cited sources and after using the function \textit{\{\textbackslash cite\{name of the cited object\}\}} the corresponding source will be added to the list of literature sources.


\textit{bibliography.bib} provides examples of some of the most commonly cited types of sources:
\begin{itemize}
    \item web pages (\textit{@online}) \cite{PvzInternetinisPuslapis},
    \item datasets (\textit{@dataset}) \cite{dataset}
    \item articles (\textit{@article}) \cite{PvzStraipsnLt, PvzStraipsnEn}, 
    \item articles from conferences (\textit{@inproceedings}) \cite{PvzKonfLt, PvzKonfEn}, 
    \item books (\textit{@book}) \cite{PvzKnygLt, PvzKnygEn}, 
    \item theses (\textit{@thesis or mastersthesis/phdthesis} \cite{PvzMagistrLt, PvzPhdEn})
    \item electronic publications (\textit{@misc}) \cite{PvzElPubLt, PvzElPubEn}
\end{itemize}

Examples are also provided for ChatGPT citation, both in general \cite{chatgpt_bendrai} and for a specific conversation \cite{chatgpt_pokalbis}.

\end{document}
